%! Function Composition and Decomposition — Unit 1, Lesson 1.3 — Student Packet Cover Sheet
% The ONE place \namedateperiod belongs (Namestrip). The packet table carries no activity and no
% exit-ticket row: this course has neither (LESSON_SHAPE.md §1, §2).
\documentclass[10pt]{article}
\usepackage{precalculus-article}
\usepackage{precalculus-boxes}
\usepackage{ltablex}
\keepXColumns

\begin{document}

% Full-bleed plum banner
\begin{tikzpicture}[remember picture, overlay]
  \fill[plum] (current page.north west)
    rectangle ([yshift=-0.9in]current page.north east);
\end{tikzpicture}
\vspace{-0.8in}

\begin{center}
  {\color{white}\textbf{\LARGE Precalculus}}\\[4pt]
  {\color{lilacmid}\large\textbf{Unit 1: Functions and Change}}\\[6pt]
  {\color{white}\textbf{\large Lesson 1.3 \quad Function Composition and Decomposition}}
\end{center}

\vspace{0.22in}
\namedateperiod
\vspace{0.18in}

\begin{learningtargetbox}
\small
\begin{enumerate}[leftmargin=1.5em, itemsep=5pt]
  \item I can name the \textit{handoff} in a composite --- the one number the inner
        function gives the outer one --- and use it to evaluate $f(g(3))$ from a
        table or from formulas.
  \item I can build $(f \circ g)(x)$ by substituting $g(x)$ for every $x$ in
        $f(x)$, and simplify the result.
  \item I can explain why $f(g(x))$ and $g(f(x))$ are usually different
        functions, and why a composite has no value when the outer function will
        not accept the handoff.
  \item I can decompose a function into an inner rule and an outer rule, and
        check my split by substituting back.
\end{enumerate}
\end{learningtargetbox}

\vspace{0.14in}

\begin{tocbox}
\small
\renewcommand{\arraystretch}{1.5}
\begin{tabularx}{\linewidth}{@{} c l X r @{}}
\rowcolor{lilac}
\textbf{\#} & \textbf{Component} & \textbf{Description} & \textbf{Score} \\
\midrule
1 & Warm-Up      & Spiral review of prerequisite skills                & \blank{1.2cm} \\
2 & Guided Notes & The handoff: find it, hand it over, take it apart   & \blank{1.2cm} \\
\rowcolor{goldbg}
3 & AP Practice  & \textbf{Extra credit} --- AP-style multiple choice and free response & $+$\,\blank{1.0cm} \\
4 & Homework     & Graded practice in new contexts \quad Due: \blank{2.2cm} & \blank{1.2cm} \\
\midrule
  & \multicolumn{2}{r}{\textbf{Total} \quad {\footnotesize (1, 2, and 4, plus any extra credit)}} & \blank{1.2cm} \\
\end{tabularx}
\end{tocbox}

\vspace{0.14in}

\begin{remindbox}
\small
Two functions can work in a row, and when they do, one number sits between them:
the \textbf{handoff}. It is the inner function's answer and the outer function's
question at the same time --- so find it first and write it down. Everything
today follows from it. Swap the two machines and you swap the handoff, which is
why $f(g(x))$ and $g(f(x))$ are usually \emph{different functions} rather than
two spellings of the same one. A handoff the outer function will not accept
means the composite has no value at all. And read backwards, the same idea takes
a function apart: name what you would compute first, and you have found the
inner rule.
\end{remindbox}

\end{document}
